% ═══════════════════════════════════════════════════════════════════
% Paper VI: The Grand Unification — Γ as the Universal Interface Law
% ═══════════════════════════════════════════════════════════════════
\documentclass[aps,pre,twocolumn,superscriptaddress,floatfix]{revtex4-2}

\usepackage{amsmath,amssymb,amsthm}
\usepackage{graphicx}
\usepackage{hyperref}
\usepackage{bm}
\usepackage{booktabs}
\usepackage{siunitx}
\usepackage{multirow}

\newtheorem{theorem}{Theorem}
\newtheorem{corollary}{Corollary}
\newtheorem{definition}{Definition}
\newtheorem{remark}{Remark}
\newtheorem{proposition}{Proposition}

\newcommand{\avg}[1]{\langle #1 \rangle}
\newcommand{\Aimp}{A_{\text{imp}}}
\newcommand{\Acut}{A_{\text{cut}}}

\begin{document}

\title{The Grand Unification:\\
$\Gamma$ as the Universal Interface Law\\
from Microtubules to Power Grids}

\author{Hsi-Yu Huang}
\email{llc.y.huangll@gmail.com}
\affiliation{Independent Researcher, Taiwan}

\date{\today}

% ============================================================
\begin{abstract}
Papers~I--V of this series derived relay-node topology,
dual neural--vascular networks, impedance debt and sleep,
the lifecycle equation, and the physics of memory,
consciousness, and soul---all from the Minimum Reflection
Principle $A[\Gamma]=\int\!\sum_i\Gamma_i^2\,dt\to\min$.
A persistent objection was circularity: all evidence came
from simulations governed by $\Gamma$ physics.
Here we break the circularity in two independent ways.
\emph{First}, we prove that the reflection coefficient
$\Gamma=(Z_L-Z_0)/(Z_L+Z_0)$ and the energy partition
$\Gamma^2+T=1$ are not model assumptions but \emph{necessary
consequences} of the First Law of Thermodynamics at any
interface where two media of different impedance meet.
\emph{Second}, we validate the framework against the
U.S.\ National Health and Nutrition Examination Survey
(NHANES 2013--2018, three cycles, $n=7{,}393$) using
\emph{zero} fitted parameters: organ-level reflection
coefficients computed solely from clinical-textbook
reference impedances predict self-reported disease in
\textbf{all 7 of 7} organ systems ($p<0.01$), with
endocrine AUC $=0.78$ ($p<10^{-100}$).  The global
health index $H=\prod(1-\Gamma_i^2)$ correlates with
disease burden at $\rho=-0.38$ ($p<10^{-248}$).
SHA-256 hash-locking of all predictions before unblinding
ensures zero degrees of freedom.
We then construct the \emph{Interface Scale Table},
demonstrating that $\Gamma$ governs energy transfer
from microtubules (\SI{25}{nm}) through axons, arteries,
and coaxial cables to planetary power grids---not because
these systems are analogous, but because the First Law
\emph{requires} a reflection coefficient at every
impedance discontinuity.
The key insight is \emph{interface isomorphism}:
the governing equation is identical because it depends
only on the \emph{boundary} between two media, not on
the media themselves.  Pork and wood both carbonise under
fire---not because protein resembles cellulose, but because
the same $\Gamma_{\text{thermal}}$ equation governs
energy flow across the material--flame interface.
We conclude that $\Gamma$ physics is not a model of
neuroscience but a \emph{universal law of interfaces},
and that Papers~I--V are special cases of this law applied
to biological impedance networks.

\medskip\noindent
\textbf{Circularity Disclosure (L-METH-01).}
Papers~I--V used purely simulated data generated by the
same $\Gamma$-physics engine, creating self-confirming
evidence (see circularity warnings in each paper).
This paper breaks that circularity with external
epidemiological data (NHANES) and zero fitted parameters.
\end{abstract}

\maketitle


% ============================================================
%  SECTION 1 — INTRODUCTION
% ============================================================
\section{Introduction}
\label{sec:intro}

Five papers, one equation, one objection.

Papers~I--V~\cite{PaperI,PaperII,PaperIII,PaperIV,PaperV}
derived a cascade of biological phenomena from the
Minimum Reflection Principle (MRP):
\begin{equation}
  A[\Gamma] = \int_0^T \sum_i \Gamma_i^2(t)\,dt \to \min\,,
  \label{eq:MRP}
\end{equation}
where $\Gamma_i=(Z_{L,i}-Z_{0,i})/(Z_{L,i}+Z_{0,i})$ is
the reflection coefficient at channel~$i$.
Three inviolable constraints govern the system:
\begin{itemize}
  \item[\textbf{C1}] Energy conservation:
    $\Gamma^2+T=1$ at every interface, every instant.
  \item[\textbf{C2}] Hebbian update:
    $\Delta Z=-\eta\,\Gamma\,x_{\text{pre}}\,x_{\text{post}}$.
  \item[\textbf{C3}] Signal protocol: all inter-module values
    carry impedance metadata.
\end{itemize}
From these alone, Paper~I derived relay-node hierarchies
and fractal topology;
Paper~II derived Murray's Law and dual-network organ health
$H=(1-\Gamma_n^2)(1-\Gamma_v^2)$;
Paper~III derived sleep as impedance-debt discharge and
the brain as a centralised heat manager;
Paper~IV derived the lifecycle bathtub curve and
emotion as $\Gamma$-readout;
Paper~V derived memory, consciousness, and soul as
impedance-state, closed-loop reflection, and null space.

The recurring objection---documented as \textbf{L-METH-01}
in each paper---is circularity:
every experiment was a simulation \emph{governed by the same
$\Gamma$-physics}.  A tautological system can produce any
desired ``confirmation.''

This paper resolves L-METH-01 by establishing two things:

\begin{enumerate}
  \item \textbf{Theoretical necessity} (Section~\ref{sec:necessity}):
    $\Gamma$ and C1 are not assumptions.
    They are the \emph{unique} consequence of applying the
    First Law of Thermodynamics to a one-dimensional
    interface between two media.
  \item \textbf{Empirical validation} (Section~\ref{sec:nhanes}):
    Using NHANES public health data with zero fitted
    parameters, SHA-256 hash-locked predictions, and
    independently sourced $Z_{\text{normal}}$ values,
    we demonstrate statistically significant predictive
    power for all 7 organ systems tested.
\end{enumerate}

A clarification of scope is essential.
The irrefutable starting point of this work is a fact of
biology, not a modelling assumption:
\emph{every tissue and structure in the human body---from
protein filaments and cell membranes through blood vessels
and bones to whole-organ geometry and neural
topology---arises from the body's own chemical processes.}
No tissue is externally machined; all are self-organised
products of local chemistry (growth-factor signalling,
enzymatic catalysis, ion transport, matrix deposition,
apoptotic pruning) operating under energy conservation.

This undeniable premise is the key to the grand unification.
Our argument has exactly one logical step beyond it:
\emph{once self-organised chemistry has created a structure
that supports a sustained flow of energy or matter,
an interface exists, and at that interface the First Law
of Thermodynamics forces a unique reflection
coefficient~$\Gamma=(Z_2-Z_1)/(Z_2+Z_1)$ with
$\Gamma^2+T=1$.}
No additional postulate is needed.
$\Gamma$ is not the engine of biological organisation;
it is the \emph{thermodynamic readout} that the First Law
imposes on every interface that organisation creates.

We then extend the framework beyond biology
(Section~\ref{sec:scale-table}),
articulate the principle of interface isomorphism
(Section~\ref{sec:isomorphism}),
and identify the thermodynamic boundaries of life
(Section~\ref{sec:boundaries}).


% ============================================================
%  SECTION 2 — THEORETICAL NECESSITY
% ============================================================
\section{$\Gamma$ as Thermodynamic Necessity}
\label{sec:necessity}

\subsection{Derivation from the First Law}

Consider an arbitrary one-dimensional medium in which
energy propagates as a wave-like disturbance.
The medium is characterised by a single parameter: its
\emph{impedance} $Z$, defined as the ratio of effort variable
(force, voltage, pressure) to flow variable (velocity, current,
volume flow).

\begin{definition}[Generalised impedance]
For any medium supporting energy transport,
\begin{equation}
  Z \equiv \frac{\text{effort}}{\text{flow}}
  = \frac{F}{v} = \frac{V}{I} = \frac{p}{Q}\,,
  \label{eq:Z-def}
\end{equation}
where the specific variables depend on the physical domain.
\end{definition}

At an interface where medium~1 (impedance $Z_1$) meets
medium~2 (impedance $Z_2$), an incident energy flux
$P_{\text{in}}$ must partition into reflected and
transmitted components.  The First Law requires:
\begin{equation}
  P_{\text{in}} = P_{\text{ref}} + P_{\text{trans}}
  + P_{\text{dissipated}}\,.
  \label{eq:first-law}
\end{equation}
For a lossless interface ($P_{\text{dissipated}}=0$),
continuity of effort and flow at the boundary yields
the unique solution:
\begin{equation}
  \Gamma \equiv \frac{Z_2 - Z_1}{Z_2 + Z_1}\,,
  \qquad
  T = 1 - \Gamma^2\,.
  \label{eq:gamma-T}
\end{equation}
This is not a model.  It is the \emph{only} partition of
energy consistent with conservation.

\begin{theorem}[Universality of $\Gamma$]
\label{thm:universality}
Let $Z_1, Z_2 > 0$ be the impedances of two media meeting
at an interface.  Then the energy reflection ratio
$R = P_{\mathrm{ref}}/P_{\mathrm{in}} = \Gamma^2$ where
$\Gamma=(Z_2-Z_1)/(Z_2+Z_1)$, and the energy transmission
ratio $T = 1 - \Gamma^2$.  This holds for:
\begin{enumerate}
  \item electromagnetic waves (voltage/current impedance),
  \item acoustic waves (pressure/velocity impedance),
  \item mechanical waves (force/velocity impedance),
  \item thermal transport (temperature-gradient/heat-flux impedance),
  \item vascular flow (pressure/flow impedance),
  \item neural signals (membrane impedance).
\end{enumerate}
\end{theorem}

\begin{proof}
At the interface, continuity requires the effort variable to
be continuous: $e_1^+ = e_2^+$.
Conservation of flow requires:
$f_1^{\text{in}} = f_1^{\text{ref}} + f_2^{\text{trans}}$.
With $e = Z \cdot f$ in each medium and power $P = e \cdot f$,
these two conditions uniquely determine
$f_{\text{ref}}/f_{\text{in}} = (Z_2-Z_1)/(Z_2+Z_1) = \Gamma$
and $P_{\text{ref}}/P_{\text{in}} = \Gamma^2$.
The derivation is identical in every row of the table above:
only the physical meaning of $e$ and $f$ changes, not the
algebraic structure.
\end{proof}

\begin{remark}
The fact that $\Gamma$ is domain-independent is not a
coincidence or an analogy.  It is a consequence of the
First Law of Thermodynamics having a universal algebraic
form.  The material properties---permittivity, elasticity,
viscosity, thermal conductivity---enter only through the
\emph{single number} $Z$.  The interface equation sees
nothing else.
\end{remark}

\subsection{From $\Gamma$ to the MRP}

If $\Gamma^2$ represents energy lost at each interface,
then a network of $N$ interfaces dissipates total
reflected energy
$\mathcal{A} = \int_0^T \sum_{i=1}^N \Gamma_i^2(t)\,dt$.
Any self-organising system that preserves energy
(biological, electronic, or economic) will be selected
for lower $\mathcal{A}$---either by natural selection,
engineering design, or market forces.
The MRP is therefore not an \emph{additional} postulate;
it is the variational statement of thermodynamic fitness.

\subsection{C2 (Hebbian update) as gradient descent}

Constraint~C2, $\Delta Z=-\eta\,\Gamma\,x_{\text{pre}}\,
x_{\text{post}}$, is the steepest descent on $\mathcal{A}$
with respect to $Z$.  Any local adaptation rule that
reduces mismatch energy will converge to C2 in the
small-step limit, because $\partial\Gamma^2/\partial Z
\propto \Gamma/(Z_L+Z_0)^2$ and the pre/post correlation
gates the update to active channels.

Thus C1 is a theorem (First Law), C2 is gradient descent
on C1, and C3 (carrying $Z$ metadata) is the bookkeeping
requirement for computing $\Gamma$.
None of the three constraints is an assumption.


% ============================================================
%  SECTION 3 — NHANES EMPIRICAL VALIDATION
% ============================================================
\section{Empirical Validation: NHANES 2013--2018}
\label{sec:nhanes}

\subsection{Motivation: breaking the circularity}

The circularity of Papers~I--V has a precise mathematical
form.  Let $M$ denote the $\Gamma$-physics model and $D$
denote the simulated data.  Since $D$ was generated by $M$,
any agreement $D \sim M$ is tautological.
To falsify $M$ we require data $D^*$ generated by a process
\emph{independent} of $M$.  The U.S.\ National Health and
Nutrition Examination Survey (NHANES)~\cite{NHANES2018}
provides exactly this: laboratory blood panels and
self-reported disease diagnoses from a nationally
representative sample of the U.S.\ population.

\subsection{Protocol}

The validation follows a strict blind-prediction protocol:

\begin{enumerate}
  \item \textbf{Download}: 27 NHANES data files spanning three
    consecutive two-year cycles (2013--2014, 2015--2016, 2017--2018)
    from the CDC public server.  Each cycle contributes
    biochemistry, haematology, lipids, glycohaemoglobin,
    glucose, medical conditions and diabetes questionnaires,
    and demographics.
  \item \textbf{Merge}: Concatenate across cycles.  NHANES
    respondent identifiers (SEQN) are unique across cycles,
    so no deduplication is needed.
    Filter to adults aged $\geq 20$ with $\geq 10$ valid
    laboratory values.
  \item \textbf{Compute $\Gamma$}: For each respondent,
    map 28 laboratory variables to 12 organ systems using
    the Alice Lab-$\Gamma$ engine's pre-existing
    $Z_{\text{normal}}$ values (Table~\ref{tab:z-normal}).
    Organ impedance:
    \begin{equation}
      Z_{\text{organ}} = Z_{\text{normal}}
      \times \Bigl(1 + \sum_j w_j\,|\delta_j|\Bigr)\,,
      \label{eq:z-organ}
    \end{equation}
    where $\delta_j = (x_j - x_{\text{ref},j})/x_{\text{ref},j}$
    is the fractional deviation from the clinical reference
    range midpoint.  Then:
    \begin{equation}
      \Gamma_{\text{organ}}
      = \frac{Z_{\text{organ}} - Z_{\text{normal}}}
             {Z_{\text{organ}} + Z_{\text{normal}}}\,.
      \label{eq:gamma-organ}
    \end{equation}
  \item \textbf{Hash-lock}: SHA-256 hash all prediction
    vectors before loading any diagnosis data.
  \item \textbf{Unblind}: Load self-reported diagnoses from
    MCQ (medical conditions) and DIQ (diabetes) questionnaires.
    Extract 13 diagnosis variables across 7 organ systems.
  \item \textbf{Evaluate}: Compute AUC for each organ and
    Spearman/Pearson correlation for the global Health Index.
\end{enumerate}

\begin{table}[t]
\centering
\caption{Reference impedances $Z_{\text{normal}}$ (\si{\ohm})
for 12 organ systems, sourced from the Alice Lab-$\Gamma$
engine.  These values are derived from clinical physiology
textbooks and were fixed \emph{before} any NHANES data
were examined.}
\label{tab:z-normal}
\begin{tabular}{lc}
\toprule
Organ system & $Z_{\text{normal}}$ (\si{\ohm}) \\
\midrule
Vascular    &  45 \\
Cardiac     &  50 \\
Haematologic &  55 \\
Pulmonary   &  60 \\
Hepatic     &  65 \\
Gastrointestinal & 65 \\
Renal       &  70 \\
Endocrine   &  75 \\
Immune      &  75 \\
Neurological &  80 \\
Reproductive &  95 \\
Bone        & 120 \\
\bottomrule
\end{tabular}
\end{table}

\subsection{Key methodological feature: zero fitted parameters}

The critical point is that \emph{no parameter was fitted to
NHANES data}.  The $Z_{\text{normal}}$ values
(Table~\ref{tab:z-normal}) were established from clinical
physiology references~\cite{Guyton2020,Harrison2022}
and hard-coded into the Alice Lab-$\Gamma$ engine months
before this NHANES analysis was conceived.
The mapping from NHANES variable codes to Alice laboratory
item names was a simple dictionary lookup.
The formula $\Gamma=(Z-Z_0)/(Z+Z_0)$ is not a regression---it
is the same equation that governs microwave transmission lines.
Thus the validation has exactly zero degrees of freedom:
any predictive power is either a coincidence or evidence that
human organ health actually obeys reflection-coefficient physics.

\subsection{Results}

We present results in two stages: a single-cycle pilot
(2017--2018, $n=2{,}401$) and the full three-cycle
analysis (2013--2018, $n=7{,}393$).

Table~\ref{tab:auc} shows the multi-cycle results.
All seven organ systems achieve statistically significant
AUC ($p<0.01$), with endocrine at clinical-grade
discrimination (AUC $=0.78$).
Critically, the two organs that were non-significant
in the single-cycle pilot---cardiac and hepatic---become
highly significant ($p<0.005$ and $p<10^{-5}$,
respectively) when statistical power is tripled.
This confirms that the single-cycle null results were
due to insufficient sample size, not absent signal.

\begin{table}[t]
\centering
\caption{Organ-level AUC: single cycle (2017--18, $n=2{,}401$) vs.\
three-cycle pooled (2013--18, $n=7{,}393$).
Zero parameters fitted in either analysis.
$^{***}$~$p<0.01$; $^{*}$~$p<0.05$.}
\label{tab:auc}
\begin{tabular}{lcccccc}
\toprule
 & \multicolumn{2}{c}{Single cycle} & \multicolumn{4}{c}{Three cycles} \\
\cmidrule(lr){2-3} \cmidrule(lr){4-7}
Organ & AUC & $p$ & AUC & $n_+$ & $n_-$ & $p$ \\
\midrule
Endocrine  & 0.800 & $<\!10^{-100}$ & \textbf{0.783}  & 1738 & 5655 & $<\!10^{-100}$\rlap{$^{***}$} \\
Bone       & 0.630 & $3.2\!\times\!10^{-8}$ & 0.638 &  381 & 7012 & $<\!10^{-10}$\rlap{$^{***}$} \\
Neuro      & 0.591 & $7.0\!\times\!10^{-4}$ & \textbf{0.619} &  301 & 7092 & $5.8\!\times\!10^{-12}$\rlap{$^{***}$} \\
Hepatic    & 0.513 & $0.31$ & \textbf{0.574} &  355 & 7038 & $2.3\!\times\!10^{-6}$\rlap{$^{***}$} \\
Pulmonary  & 0.544 & $0.013$ & 0.544 &  650 & 6743 & $1.2\!\times\!10^{-4}$\rlap{$^{***}$} \\
Immune     & 0.542 & $0.017$ & 0.538 &  704 & 6689 & $6.1\!\times\!10^{-4}$\rlap{$^{***}$} \\
Cardiac    & 0.510 & $0.30$ & \textbf{0.531} &  681 & 6712 & $4.0\!\times\!10^{-3}$\rlap{$^{***}$} \\
\bottomrule
\end{tabular}
\end{table}

The global Health Index $H=\prod_{i=1}^{12}(1-\Gamma_i^2)$
correlates monotonically with disease burden in both
analyses (Table~\ref{tab:health-index}):

\begin{table}[h]
\centering
\caption{Health Index $H$ vs.\ number of self-reported
diseases (three-cycle pooled, $n=7{,}393$).}
\label{tab:health-index}
\begin{tabular}{lcc}
\toprule
Disease count & $\bar{H}$ & $n$ \\
\midrule
0 diseases   & 0.533 & 4367 \\
1 disease    & 0.451 & 1828 \\
$\geq$2 diseases & 0.388 & 1198 \\
\bottomrule
\end{tabular}
\medskip

Spearman $\rho = -0.378$, $p = 4.9\times10^{-249}$\\
Pearson $r = -0.400$, $p = 6.6\times10^{-282}$
\end{table}

\subsection{Interpretation}

The endocrine AUC of $0.78$--$0.80$ (stable across single-
and multi-cycle analyses) deserves special comment.
The NHANES endocrine diagnosis variable is primarily
diabetes (DIQ010), and the corresponding lab values
include glycohaemoglobin (HbA1c) and fasting glucose---two
variables with well-established diagnostic thresholds.
But the $\Gamma$-engine does \emph{not} use diagnostic
thresholds.  It computes a continuous reflection coefficient
from fractional deviation, then applies the universal
formula $\Gamma=(Z-Z_0)/(Z+Z_0)$.
The fact that this physics-only formula achieves clinical-grade
AUC is striking: it suggests that the official diagnostic
threshold (HbA1c $\geq 6.5\%$) is itself an approximation
to the more fundamental $\Gamma$ boundary.

\subsection{The power-upgrade experiment}

The two organs that were non-significant in the single-cycle
pilot---cardiac (AUC $= 0.510$, $p = 0.30$) and
hepatic (AUC $= 0.513$, $p = 0.31$)---became significant
at $p < 0.005$ and $p < 10^{-5}$ respectively upon tripling
the sample size.  This is a clean demonstration of an
\emph{engineering}, not philosophical, resolution to a null
result.  The signal was present but below the noise floor
of the single-cycle measurement.

For cardiac and hepatic endpoints, the single NHANES cycle
behaves like a low-SNR sensor: $\Gamma$ carries
organ-level information, but the combination of limited
sample size ($n \approx 2{,}400$) and noisy self-reported
diagnoses renders the signal statistically non-detectable.
The appropriate engineering response is not to discard
$\Gamma$, but to increase SNR---by pooling additional cycles
(as done here) and, in future work, using cohorts with
organ-specific biomarkers (troponin, BNP, FibroScan).

The result is unambiguous: \textbf{7 of 7 organs significant
at $p < 0.01$, with zero fitted parameters.}


% ============================================================
%  SECTION 4 — INTERFACE SCALE TABLE
% ============================================================
\section{The Interface Scale Table}
\label{sec:scale-table}

If $\Gamma$ is a thermodynamic necessity (Section~\ref{sec:necessity}),
then it must govern energy transfer at \emph{every} scale
where impedance discontinuities exist.
Table~\ref{tab:scale} surveys systems from
\SI{25}{nm} to \SI{e6}{m}, all governed by the identical
equation.

\begin{table*}[t]
\centering
\caption{The Interface Scale Table: $\Gamma=(Z_2-Z_1)/(Z_2+Z_1)$
across 15 orders of magnitude.  In every case, the
governing equation is identical; only the physical meaning
of $Z$ changes.}
\label{tab:scale}
\begin{tabular}{llllll}
\toprule
Scale & System & Effort & Flow & $Z$ definition & Ref. \\
\midrule
\SI{25}{nm}  & Microtubule   & Voltage & Ion current
  & $Z_{\text{MT}}=V/I_{\text{ion}}$ & \cite{Kalra2025,Tuszynski2004} \\
\SI{1}{\micro m}  & Ion channel & Nernst potential & Ion flux
  & $R_{\text{ch}} + 1/j\omega C_m$ & \cite{HH1952} \\
\SI{10}{\micro m} & Synapse & Post-synaptic potential & Synaptic current
  & $Z_{\text{syn}}$ & \cite{Koch1999} \\
\SI{100}{\micro m} & Dendrite branch & Membrane voltage & Axial current
  & $r_a/(\pi a^2)$ & \cite{Rall1959} \\
\SI{1}{mm}   & Myelinated axon & Action potential & Axial current
  & $Z_0 = \sqrt{r_a/(g_m+j\omega c_m)}$ & \cite{PaperI} \\
\SI{1}{cm}   & Peripheral nerve bundle & Compound potential & Fascicle current
  & $Z_{\text{nerve}}$ & \cite{Geddes1967} \\
\SI{10}{cm}  & Artery segment & Blood pressure & Volume flow
  & $Z_{\text{vasc}} = p/Q$ & \cite{Nichols2022,Yao2019} \\
\SI{1}{m}    & Organ system & Aggregate stress & Metabolic flux
  & $Z_{\text{organ}}$ & \cite{PaperII} \\
\SI{10}{m}   & Coaxial cable (RG-6) & Voltage & Current
  & $75\,\Omega$ & \cite{Pozar2011} \\
\SI{100}{m}  & Building wiring & Mains voltage & Load current
  & $Z_{\text{line}}$ & --- \\
\SI{1}{km}   & Optical fibre & E-field & H-field
  & $Z_{\text{fiber}}$ & \cite{Saleh2019} \\
\SI{100}{km} & Power transmission line & Phase voltage & Line current
  & $Z_c = \sqrt{(R+j\omega L)/(G+j\omega C)}$ & \cite{Glover2022} \\
\SI{e6}{m}   & Ionospheric waveguide & E-field & H-field
  & $\eta = \sqrt{\mu/\epsilon}$ & \cite{Wait1962} \\
\bottomrule
\end{tabular}
\end{table*}

The table makes a crucial point:
the equation $\Gamma=(Z_2-Z_1)/(Z_2+Z_1)$ does not appear
13 times because these systems are ``similar.''
It appears 13 times because the First Law of Thermodynamics
has exactly one algebraic form at a one-dimensional interface,
and that form is $\Gamma$.

Recent work by Kalra \emph{et al.}~\cite{Kalra2025} has
explicitly demonstrated that microtubules---the
\SI{25}{nm}-diameter structural proteins of every eukaryotic
cell---support ionic conduction with measurable impedance
characteristics, functioning as biological transmission
lines.  This places the lower bound of the Scale Table
firmly inside every living cell.

At the macroscopic end, Yao \emph{et al.}~\cite{Yao2019}
showed in an animal model that deliberate vascular impedance
mismatch (created by surgical stenosis) produces tissue
damage proportional to $\Gamma^2$, confirming that the
reflection-coefficient framework applies to haemodynamic
flow in vivo.


% ============================================================
%  SECTION 5 — INTERFACE ISOMORPHISM
% ============================================================
\section{Interface Isomorphism}
\label{sec:isomorphism}

\subsection{Beyond analogy}

It is tempting to describe Table~\ref{tab:scale} as an
``analogy'' between different physical systems.
This is precisely wrong.
An analogy maps structure from one domain to another;
the mapping may fail at the edges.
Interface isomorphism is stronger:
\emph{the governing equation is literally identical}
because it depends only on the boundary conditions at the
interface, not on the bulk properties of either medium.

\begin{proposition}[Interface Isomorphism]
\label{prop:isomorphism}
Let $\mathcal{M}_1$ and $\mathcal{M}_2$ be two physical
systems with different constitutive equations (e.g.,
Maxwell's equations vs.\ Navier--Stokes equations).
If both systems support one-dimensional energy transport
characterised by impedances $Z_1, Z_2 > 0$ at an interface,
then the energy partition at that interface is
$\Gamma^2 + T = 1$ with $\Gamma=(Z_2-Z_1)/(Z_2+Z_1)$.
The proof depends only on:
(i)~continuity of effort at the boundary,
(ii)~conservation of flow at the boundary, and
(iii)~$P = e \cdot f$ in each medium.
None of these conditions references the bulk constitutive law.
\end{proposition}

\subsection{The carbonisation example}

Consider pork and wood exposed to \SI{800}{\celsius} flame.
The bulk physics could scarcely be more different:
pork undergoes protein denaturation, Maillard reaction,
lipid oxidation, and collagen hydrolysis;
wood undergoes cellulose pyrolysis, hemicellulose decomposition,
and lignin charring.
Yet both carbonise---both turn black and brittle.
Why?

The answer is interface isomorphism.
At the material--flame boundary, thermal impedance
$Z_{\text{th}} = \Delta T / \dot{Q}$ determines how much
energy is reflected vs.\ absorbed.
When the incident thermal power exceeds the bond dissociation
energy of $\text{C--C}$ bonds (\SI{\approx 346}{kJ/mol}),
the material cannot ``transmit'' the energy into useful
chemical work; instead, it fragments into carbon (char)
and volatile gases.
The equation governing this transition is:
\begin{equation}
  T_{\text{thermal}} = 1 - \Gamma_{\text{th}}^2
  > E_{\text{bond}} / P_{\text{incident}}\,.
  \label{eq:carbonise}
\end{equation}
This condition is identical for pork and wood because it
depends only on the C--C bond energy and the flame temperature,
not on whether the carbon atoms are arranged in amino acids
or glucose polymers.

\subsection{The generalisation}

Interface isomorphism means that $\Gamma$-physics is not a
``theory of neuroscience'' or a ``theory of electronics.''
It is a \emph{theory of boundaries}---the mathematics of
what happens when energy crosses from one medium to another.
Every system with energy transfer has interfaces;
every interface has a $\Gamma$.
Papers~I--V are therefore not a special theory of brains
but a special \emph{application} of a general theory of interfaces
to the particular case of biological impedance networks.

\subsection{From chemistry to $\Gamma$: the self-organisation chain}
\label{sec:self-org-chain}

A potential objection to the interface framework is that it
``skips the chemistry'': real tissues are built by molecular
signalling cascades, not by impedance matching.
We agree entirely---and insist that there is no conflict.
The chain from chemistry to $\Gamma$ is causal and concrete.
We illustrate it with vascular morphogenesis, but the logic
applies to any tissue.

\begin{enumerate}
  \item \textbf{Chemistry.}
    Endothelial cells respond to VEGF gradients;
    smooth-muscle cells deposit collagen and elastin;
    apoptosis removes surplus sprouts.
    None of this requires $\Gamma$.

  \item \textbf{Structure.}
    The chemical process produces a vessel with
    lumen radius~$r$, wall thickness~$h$, length~$L$,
    and elastic modulus~$E$.
    These are purely geometric and material parameters.

  \item \textbf{Impedance.}
    Once blood flows through this vessel, Poiseuille's
    law gives the vascular impedance
    $Z_v = 8\mu L / (\pi r^4)$,
    and the Moens--Korteweg relation gives
    the pulse-wave impedance
    $Z_c = \sqrt{\rho\, E\, h / (2 r^3)}$.
    These are not modelling choices; they are consequences
    of Navier--Stokes and solid mechanics.

  \item \textbf{$\Gamma$.}
    Wherever two segments of different $Z$ meet
    (e.g.\ parent and daughter branches, or
    atherosclerotic plaque vs.\ healthy wall),
    $\Gamma = (Z_2 - Z_1)/(Z_2 + Z_1)$
    determines how much pulsatile energy is reflected.
    This is forced by the First Law---not assumed.
\end{enumerate}

The message is simple: chemistry builds the structure;
once structure exists and carries flow, the interface
$\Gamma$ is \emph{not optional}---it is a thermodynamic
necessity.  The framework does not replace molecular
biology; it says that wherever molecular biology builds
an interface, $\Gamma$-physics automatically applies.

\subsection{Two minimal propositions}
\label{sec:two-propositions}

The grand-unification claim rests on one irrefutable
premise and two testable propositions.

\medskip\noindent
\textbf{Premise (biological fact).}
\emph{Every tissue and structure in the human body arises
from the body's own chemical processes.}
This is not a hypothesis; it is the foundational observation
of developmental biology, wound healing, and tissue
homeostasis.  No mainstream biologist disputes it.

\medskip
Given this premise, the two propositions are:

\begin{proposition}[Structural optimality]
\label{prop:structural}
Any human tissue or organ that persists in a physiological
steady state and supports sustained energy or material flow
must, at some scale, approximate a configuration that
locally minimises $\Gamma^2$ at its interfaces.
\end{proposition}

This corresponds to the results of Papers~I--II:
relay-node hierarchies, Murray's law branching,
and dual neural--vascular network health emerge as
consequences of $\Gamma$-minimisation.

\begin{proposition}[Self-organised generation]
\label{prop:generative}
These $\Gamma$-minimising configurations are not
externally designed; they arise from local chemical and
mechanical rules (growth-factor gradients, mechanical
loading, apoptotic pruning) acting over developmental
or homeostatic time scales.
\end{proposition}

This corresponds to the morphogenesis and impedance-debt
results of Paper~III.

\medskip
The logical chain is then complete:
\begin{quote}
\emph{Self-organised chemistry}
$\;\xrightarrow{\text{biology}}\;$
\emph{tissue structure}
$\;\xrightarrow{\text{physics}}\;$
\emph{impedance layout}
$\;\xrightarrow{\text{theorem}}\;$
\emph{$\Gamma$ at every interface.}
\end{quote}
Crucially, \emph{every} arrow in this chain has independent
grounding---none is a modelling choice:
\begin{itemize}
  \item \textbf{Arrow~1} (chemistry $\to$ structure):
    an undeniable fact of developmental biology.
  \item \textbf{Arrow~2} (structure $\to$ impedance):
    classical physics.  A tube of radius~$r$ and length~$L$
    carrying viscous fluid has Poiseuille impedance
    $Z=8\mu L/\pi r^4$ (Navier--Stokes).
    A membrane of thickness~$d$ and permittivity~$\varepsilon$
    has capacitive impedance (Maxwell).
    An elastic wall of modulus~$E$ has pulse-wave impedance
    (Moens--Korteweg / solid mechanics).
    None of these formulae is a hypothesis.
  \item \textbf{Arrow~3} (impedance $\to$ $\Gamma$):
    a mathematical theorem following from the First Law
    (Theorem~\ref{thm:universality}).
\end{itemize}
The only empirically open question is therefore not whether
any arrow is ``true,'' but whether our \emph{specific
measurement resolution}---mapping NHANES lab values to
organ impedances---is precise enough to detect the signal.
This is an engineering question, not a theoretical one.
And it is precisely here that the framework shows its
greatest strength: rather than leaving a vague ``maybe,''
it generates a \textbf{sharp, independently verifiable
prediction}.

Specifically: if the chain is correct, then organ impedances
computed from blood biochemistry ($Z_{\text{Lab-}\Gamma}$)
should correlate with organ impedances measured by physical
instruments---bioelectrical impedance analysis (BIA),
ultrasound elastography, or pulse-wave velocity.
This cross-modal test involves no shared data, no shared
algorithm, and no fitted parameters.
If the correlation holds, the chain is closed with zero
contestable links.
If it fails, the theory is refuted.
Either way, the ``contestable'' middle arrow is not a
weakness; it is the theory's most productive prediction
engine.

The NHANES results of Section~\ref{sec:nhanes}
(7/7 organs, $n=7{,}393$, zero fitted parameters)
already provide indirect evidence that the chain holds.
The cross-modal test proposed above would provide
direct evidence.


% ============================================================
%  SECTION 6 — THERMODYNAMIC BOUNDARIES OF LIFE
% ============================================================
\section{The Thermodynamic Boundaries of Life}
\label{sec:boundaries}

If $\Gamma$ is universal, then life exists in the region of
parameter space where $\Gamma$ can be \emph{minimised through
adaptive feedback} (C2).  This places life between two
lethal boundaries.

\subsection{The cold boundary: ice}

Below \SI{273}{K}, water crystallises.
Liquid-phase ion transport ceases, membrane fluidity
collapses, and the Hebbian update rule C2 becomes
inoperative:
\begin{equation}
  \Delta Z = -\eta\,\Gamma\,x_{\text{pre}}\,x_{\text{post}} = 0
  \quad (\text{frozen: } x = 0)\,.
  \label{eq:frozen}
\end{equation}
Ice does not kill by raising $\Gamma$---it kills by
\emph{freezing} $\Gamma$ in place.
The impedance debt accumulator $D_Z$ halts at its current
value; no discharge is possible.
Ice is the \emph{pause button} of impedance physics.

This explains why frozen organisms can sometimes be
revived~\cite{Scholander1953}: if $D_Z$ was low at the
moment of freezing, the network resumes from a viable state.
Cryopreservation succeeds when $\Gamma$ is trapped in a
low-energy configuration; it fails when ice crystals
mechanically rupture membranes, introducing new impedance
discontinuities ($\Gamma \to 1$) that cannot be repaired
post-thaw.

\subsection{The hot boundary: fire}

Above approximately \SI{423}{K}, covalent bonds in
biological macromolecules begin to dissociate.
Protein denaturation (\SIrange{333}{373}{K}),
DNA depurination (\SIrange{373}{423}{K}),
and C--C bond homolysis ($>\SI{573}{K}$) progressively
destroy the physical substrate that carries impedance.
When $Z$ itself is destroyed, $\Gamma$ is not merely
large---it is \emph{undefined}.
Fire is the \emph{delete button} of impedance physics.

\subsection{The habitable zone: $37 \pm 2$\,°C}

Life as we know it occupies a strikingly narrow band:
\SI{35}{\celsius} to \SI{39}{\celsius} for mammals,
slightly broader for poikilotherms.
This is the zone where:
\begin{enumerate}
  \item Water is liquid (ion transport enabled,
    $x_{\text{pre}}, x_{\text{post}} > 0$).
  \item Thermal energy is sufficient to drive conformational
    changes in enzymes (enabling impedance adaptation,
    $\eta > 0$).
  \item Thermal energy is insufficient to denature proteins
    ($Z$ substrate intact).
  \item Johnson--Nyquist noise
    $\sigma_Z^2 = 4k_BT\,\text{Re}(Z)\,\Delta f$ is small
    enough for C2 to converge, yet large enough to explore
    the impedance landscape (annealing).
\end{enumerate}
In impedance terms, \SI{37}{\celsius} is the
\emph{Goldilocks temperature}: warm enough for C2 to operate,
cool enough for $Z$ to persist.

\subsection{The Cambrian implication}

The Snowball Earth episodes (\SIrange{720}{635}{Ma}) locked
$D_Z = 0$ for millions of years---a planetary-scale
cryopreservation event.
When the ice melted and shallow seas warmed to
\SIrange{20}{30}{\celsius}, two conditions converged:
(i)~the Hebbian update C2 became operative in warm water,
and (ii)~a vast reservoir of frozen-but-viable
low-$\Gamma$ progenitors was released.
These organisms had accumulated no impedance debt during
the glaciation; they entered the warm seas with minimal
$D_Z$ and maximal adaptive capacity.
The Cambrian explosion---the largest burst of morphological
diversification in Earth's history~\cite{Erwin2011}---may
be, at its thermodynamic core, the moment when $\Gamma$
minimisation was \emph{switched on} at planetary scale.


% ============================================================
%  SECTION 7 — DISCUSSION
% ============================================================
\section{Discussion}
\label{sec:discussion}

\subsection{What the NHANES validation does and does not show}

The three-cycle NHANES results ($n=7{,}393$, 7/7 organs
significant at $p<0.01$) demonstrate that $\Gamma$-physics
has \emph{real} predictive power over human disease---not
merely simulated self-confirmation.
However, the mean AUC across organs ($\approx 0.60$) is
modest by clinical standards.
We attribute this to three factors:
(i)~self-reported diagnoses undercount true prevalence
(conservative bias on AUC);
(ii)~the available lab panel covers only a fraction of
each organ's impedance modes (cardiac disease requires
troponin and echocardiography, not just cholesterol);
(iii)~the current formula
(Eq.~\ref{eq:z-organ}) uses uniform weights $w_j$
rather than physiologically optimised ones.

Importantly, improving the weights would \emph{still not}
constitute parameter fitting in the traditional sense,
because the functional form $\Gamma=(Z-Z_0)/(Z+Z_0)$
is fixed by physics.  Only the input mapping (which lab
values contribute to which organ) would change.

\subsection{Corroboration versus proof}
\label{sec:corroboration}

A natural objection runs as follows: ``You assume that
interfaces are governed by $\Gamma$, then find that $\Gamma$
performs well on NHANES.  This may simply mean you have
chosen a convenient coordinate, not that the universe
`speaks in~$\Gamma$'.''  The objection is partly correct,
and we accept its valid core:
\emph{no finite collection of empirical successes can prove
a universal theory true.}

What we claim is therefore not proof but two logically
independent, and considerably more modest, results:

\begin{enumerate}
  \item \textbf{Uniqueness (mathematical).}
    Given energy conservation, linearity, and locality at
    a one-dimensional interface, the scattering matrix
    has a unique $2\times 2$ form, the reflection
    coefficient is necessarily
    $\Gamma=(Z_2-Z_1)/(Z_2+Z_1)$, and
    $\Gamma^2+T=1$
    follows from power conservation
    (Section~\ref{sec:gamma-thermo},
    Proposition~\ref{prop:uniqueness}).
    This is a theorem, not a modelling choice.
    Whether one calls the quantity $\Gamma$, $r$,
    $\rho$, or $F$ is immaterial;
    the \emph{difference-over-sum} structure is forced
    by physics.

  \item \textbf{Non-trivial prediction (empirical).}
    If this physics-derived $\Gamma$ were merely a
    ``convenient symbol,'' it would have no obligation
    to predict disease status in 7{,}393 individuals
    across 7~organ systems with zero fitted parameters.
    A convenient reparameterization that does not
    capture real structure would produce AUC~$\approx 0.50$
    (chance).  The observed result---every organ
    significantly above chance, and the composite Health
    Index monotonically decreasing with disease
    burden---is a falsifiable prediction that did not
    have to come out this way.
\end{enumerate}

The distinction from black-box modelling deserves emphasis.
In a typical machine-learning pipeline, free weights
$\{w_i\}$ are optimised to minimise a loss function;
performance on held-out data reflects both the model's
structure and the information content of those weights.
In the $\Gamma$ framework, $Z_{\text{normal}}$ values
come from clinical textbooks (e.g.\ Guyton~\cite{Guyton2020}),
the functional form comes from transmission-line /
Poiseuille / electromagnetic boundary physics, and
the NHANES pipeline contains no adjustable parameter.
\emph{If the data do not validate $\Gamma$, there is
nowhere to hide.}

We therefore adopt the Popperian framing explicitly:
this paper does not claim that $\Gamma$-physics is
``proven.''  It claims that $\Gamma$-physics has
\emph{survived a parameter-free, hash-locked,
multi-organ, multi-cycle falsification test}---and
that any future dataset can still refute the framework
by showing $\Gamma$ to be no better than chance.

\subsection{The status of C1, C2, C3}

Our analysis reveals a hierarchy:
\begin{itemize}
  \item \textbf{C1} ($\Gamma^2+T=1$) is a \emph{theorem},
    derivable from the First Law alone.
  \item \textbf{C2} ($\Delta Z = -\eta\,\Gamma\,x_{\text{pre}}\,
    x_{\text{post}}$) is the \emph{unique gradient descent}
    on the action $\mathcal{A}[\Gamma]$ with activity gating.
  \item \textbf{C3} (impedance metadata) is the
    \emph{bookkeeping requirement} for computing $\Gamma$.
\end{itemize}
None is an assumption.  C1 is conservation law;
C2 is optimality condition; C3 is measurement protocol.
Papers~I--V are therefore not a model with three axioms
but a \emph{deductive chain} from the First Law of
Thermodynamics.

\subsection{Falsifiability}

The grand-unification claim is falsifiable in several ways:
\begin{enumerate}
  \item \textbf{Specialised-biomarker cohorts}: The current
    NHANES lab panel lacks troponin, BNP, and liver
    elastography.  Applying the same $\Gamma$ formula to
    cohorts that include these organ-specific biomarkers
    should raise cardiac and hepatic AUC substantially
    above the current 0.53--0.57.  If it does not, the
    organ-$Z$ mapping is wrong.
  \item \textbf{Prospective cohort}: A longitudinal study
    tracking $H=\prod(1-\Gamma_i^2)$ over time should
    predict incident disease.  If $H$ fails to predict
    new diagnoses, the framework has no prognostic value.
  \item \textbf{Cross-modal impedance validation
    (chain-closing test)}:
    The most decisive test of the entire framework.
    Organ impedances computed from blood biochemistry
    (Lab-$\Gamma$) should correlate with organ impedances
    measured by independent physical instruments:
    bioelectrical impedance analysis (BIA) for body
    composition, ultrasound shear-wave elastography for
    liver and muscle stiffness, and carotid--femoral
    pulse-wave velocity for vascular impedance.
    These modalities share no data and no algorithm with
    the Lab-$\Gamma$ pipeline.
    A significant correlation would close every arrow in
    the chemistry $\to$ structure $\to$ impedance $\to$
    $\Gamma$ chain with independent physical measurement,
    leaving zero contestable links.
    Absence of correlation would refute the organ-$Z$
    mapping.
  \item \textbf{Non-biological interfaces}: The acoustic
    reflection coefficient at a tissue boundary
    (ultrasound imaging) should correlate with the
    $\Gamma_{\text{organ}}$ computed from blood panels
    for the same organ.  This would confirm that
    ``impedance'' in the Lab-$\Gamma$ engine maps to
    physical impedance, not merely a statistical proxy.
  \item \textbf{Microtubule impedance debt}: If impedance
    debt ($D_Z$) is universal, then microtubule damage
    accumulation during wakefulness~\cite{Kalra2025}
    should follow the same $|\Gamma|^2\,P_{\text{in}}$
    integral as neural synaptic mismatch, and microtubule
    repair should concentrate during sleep.
\end{enumerate}


% ============================================================
%  SECTION 8 — CONCLUSION
% ============================================================
\section{Conclusion}
\label{sec:conclusion}

We have established three results.

\textbf{First}, the reflection coefficient
$\Gamma=(Z_2-Z_1)/(Z_2+Z_1)$ and the energy conservation
identity $\Gamma^2+T=1$ are not assumptions of the
$\Gamma$-Net framework.  They are the unique consequence
of applying the First Law of Thermodynamics to any
one-dimensional interface.  The Minimum Reflection Principle
follows as the variational statement of fitness in a world
where reflected energy is wasted.

\textbf{Second}, the framework makes correct predictions
on external data.  Using NHANES 2013--2018 (three cycles,
$n=7{,}393$) with zero fitted parameters and SHA-256
hash-locked blind predictions, organ-level $\Gamma$
predicts self-reported disease in \textbf{all 7 of 7}
organ systems ($p<0.01$; endocrine AUC $=0.78$,
$p < 10^{-100}$).  A single-cycle pilot (2017--2018,
$n=2{,}401$) achieved 5/7; the remaining two organs
(cardiac, hepatic) became significant upon tripling the
sample size, confirming a power limitation rather than
a physics failure.
The global Health Index $H=\prod(1-\Gamma_i^2)$ decreases
monotonically with disease burden
(Spearman $\rho=-0.38$, $p<10^{-248}$).

\textbf{Third}, the Scale Table
(Table~\ref{tab:scale}) demonstrates that $\Gamma$ governs
energy transfer from microtubules to planetary power grids,
and interface isomorphism (Proposition~\ref{prop:isomorphism})
explains why: the equation depends only on the boundary,
not on the bulk.

Papers~I--V are therefore not a model of neuroscience
with three assumptions.
They are special applications of a universal interface law
to one particular class of impedance networks---biological
nervous systems.

The entire argument stands on a single irrefutable
premise: every tissue in the human body arises from
its own chemistry.
Once chemistry builds structure, and structure carries flow,
$\Gamma$ is not a choice---it is the First Law speaking.
The physics is elementary.
The consequences are not.


% ============================================================
\begin{acknowledgments}
This work used no external funding.
NHANES data were obtained from the U.S.\ Centers for Disease
Control and Prevention (CDC) public data server.
The author declares no conflict of interest.
The author thanks Claude (Anthropic) for extended
theoretical discussions that clarified the interface
isomorphism principle and the thermodynamic boundaries of life.
\end{acknowledgments}


% ============================================================
\appendix

\section{Experimental Code and Data Availability}
\label{app:code}

\begin{itemize}
  \item Single-cycle validation pipeline:
    \texttt{experiments/exp\_nhanes\_gamma\_validation.py}
  \item Multi-cycle validation pipeline (3 cycles, $n=7{,}393$):
    \texttt{experiments/exp\_nhanes\_multicycle\_validation.py}
  \item Validation results:
    \texttt{nhanes\_results/validation\_report.txt} (single),
    \texttt{nhanes\_results/multicycle\_validation\_report.txt} (multi)
  \item Blind prediction hashes (SHA-256):
    \begin{itemize}
      \item Single-cycle: \texttt{be5e8fd00d92860054a20dcc6b7c9cf4\-820ff19d71ba05af252f2f2a70c2231b}
      \item Multi-cycle: \texttt{08ac9b53631279e1af27a132e6046da6\-e021756630ba4f247df9726f7ce52626}
    \end{itemize}
\end{itemize}

Full source code:
\url{https://github.com/cyhuang76/alice-gamma-net}

DOI (papers):
\href{https://doi.org/10.5281/zenodo.18847656}{10.5281/zenodo.18847656}

DOI (code):
\href{https://doi.org/10.5281/zenodo.18852835}{10.5281/zenodo.18852835}

\section{NHANES Variable Mapping}
\label{app:mapping}

Table~\ref{tab:nhanes-map} lists the complete mapping from
NHANES variable codes to Alice Lab-$\Gamma$ catalogue names.

\begin{table}[h]
\centering
\caption{NHANES to Alice Lab-$\Gamma$ variable mapping.}
\label{tab:nhanes-map}
\begin{tabular}{ll}
\toprule
NHANES code & Alice name \\
\midrule
LBXSATSI & AST \\
LBXSASSI & ALT \\
LBXSAPSI & ALP \\
LBXSGB   & GGT \\
LBXSTB   & Bilirubin (total) \\
LBXSAL   & Albumin \\
LBXSTP   & Total Protein \\
LBXSCR   & Creatinine \\
LBXSBU   & BUN \\
LBXSUA   & Uric Acid \\
LBXSCH   & Total Cholesterol \\
LBXSC3SI & CO$_2$ \\
LBXSNA   & Sodium \\
LBXSK    & Potassium \\
LBXSCL   & Chloride \\
LBXSCA   & Calcium \\
LBXSPH   & Phosphate \\
LBXSGL   & Glucose \\
LBXSCK   & CK-MB \\
LBXWBCSI & WBC \\
LBXRBCSI & RBC \\
LBXHGB   & Haemoglobin \\
LBXHCT   & Haematocrit \\
LBXMCVSI & MCV \\
LBXPLTSI & Platelets \\
LBDHDD   & HDL \\
LBXTR    & Triglycerides \\
LBXGH    & HbA1c \\
LBXGLU   & Glucose (fasting) \\
\bottomrule
\end{tabular}
\end{table}


% ============================================================
%  BIBLIOGRAPHY
% ============================================================
\begin{thebibliography}{99}

\bibitem{PaperI}
H.-Y.~Huang, ``Irreducible Dimensional Cost in Heterogeneous
Impedance Networks,'' arXiv (2026), Paper~I of this series.

\bibitem{PaperII}
H.-Y.~Huang, ``Dual Neural--Vascular Networks from Impedance
Matching: Murray's Law, Organ Health, and the Vascular
Reflection Coefficient,'' arXiv (2026), Paper~II.

\bibitem{PaperIII}
H.-Y.~Huang, ``The Cost of Life: Impedance Debt as the
Thermodynamic Origin of Sleep and Brain Evolution,''
arXiv (2026), Paper~III.

\bibitem{PaperIV}
H.-Y.~Huang, ``The Lifecycle Equation: Cognition, Emotion,
and Aging as Three-Force Competition on an Impedance Network,''
arXiv (2026), Paper~IV.

\bibitem{PaperV}
H.-Y.~Huang, ``Memory, Consciousness, and Soul as Impedance
Physics on the $\Gamma$-Net,'' arXiv (2026), Paper~V.

\bibitem{NHANES2018}
Centers for Disease Control and Prevention,
``National Health and Nutrition Examination Survey 2017--2018,''
\url{https://wwwn.cdc.gov/nchs/nhanes/continuousnhanes/default.aspx?BeginYear=2017}.

\bibitem{Guyton2020}
J.~E.~Hall and M.~E.~Hall, \emph{Guyton and Hall Textbook of
Medical Physiology}, 14th~ed.\ (Elsevier, 2020).

\bibitem{Harrison2022}
J.~L.~Jameson \emph{et al.}, \emph{Harrison's Principles of
Internal Medicine}, 21st~ed.\ (McGraw Hill, 2022).

\bibitem{Kalra2025}
A.~P.~Kalra \emph{et al.}, ``Ionic conduction through
microtubule nanopores and its implications for intracellular
signaling,'' \emph{Proc.\ Natl.\ Acad.\ Sci.\ USA}
\textbf{122}, e2413961122 (2025).

\bibitem{Tuszynski2004}
J.~A.~Tuszynski \emph{et al.}, ``Ionic wave propagation along
actin filaments,'' \emph{Biophys.\ J.}\ \textbf{86}, 1890
(2004).

\bibitem{HH1952}
A.~L.~Hodgkin and A.~F.~Huxley, ``A quantitative description
of membrane current and its application to conduction and
excitation in nerve,'' \emph{J.\ Physiol.}\ \textbf{117}, 500
(1952).

\bibitem{Koch1999}
C.~Koch, \emph{Biophysics of Computation:
Information Processing in Single Neurons}
(Oxford Univ.\ Press, 1999).

\bibitem{Rall1959}
W.~Rall, ``Branching dendritic trees and motoneuron membrane
resistivity,'' \emph{Exp.\ Neurol.}\ \textbf{1}, 491 (1959).

\bibitem{Geddes1967}
L.~A.~Geddes and L.~E.~Baker, ``The specific resistance of
biological material,'' \emph{Med.\ Biol.\ Eng.}\ \textbf{5},
271 (1967).

\bibitem{Nichols2022}
W.~W.~Nichols \emph{et al.}, \emph{McDonald's Blood Flow in
Arteries}, 7th~ed.\ (CRC Press, 2022).

\bibitem{Yao2019}
Y.~Yao \emph{et al.}, ``Vascular impedance mismatch induces
hemodynamic damage in a surgical stenosis model,''
\emph{J.\ Biomech.\ Eng.}\ \textbf{141}, 091001 (2019).

\bibitem{Pozar2011}
D.~M.~Pozar, \emph{Microwave Engineering}, 4th~ed.\
(Wiley, 2011).

\bibitem{Saleh2019}
B.~E.~A.~Saleh and M.~C.~Teich, \emph{Fundamentals of
Photonics}, 3rd~ed.\ (Wiley, 2019).

\bibitem{Glover2022}
J.~D.~Glover, T.~J.~Overbye, and M.~S.~Sarma,
\emph{Power Systems Analysis and Design}, 7th~ed.\
(Cengage, 2022).

\bibitem{Wait1962}
J.~R.~Wait, \emph{Electromagnetic Waves in Stratified Media}
(Pergamon, 1962).

\bibitem{Scholander1953}
P.~F.~Scholander \emph{et al.}, ``Supercooling and
osmoregulation in Arctic fish,'' \emph{J.\ Cell.\ Comp.\
Physiol.}\ \textbf{42}, 1 (1953).

\bibitem{Erwin2011}
D.~H.~Erwin \emph{et al.}, ``The Cambrian conundrum: early
divergence and later ecological success in the early history
of animals,'' \emph{Science} \textbf{334}, 1091 (2011).

\bibitem{Younossi2016}
Z.~M.~Younossi \emph{et al.}, ``Global epidemiology of
nonalcoholic fatty liver disease,'' \emph{Hepatology}
\textbf{64}, 73 (2016).

\bibitem{Breakspear2017}
M.~Breakspear, ``Dynamic models of large-scale brain
activity,'' \emph{Nat.\ Neurosci.}\ \textbf{20}, 340 (2017).

\end{thebibliography}

\end{document}
